% Propietario conservado: /Users/ruben/Documents/New project/output/EXTREMA_MEDIA_RAZON_RECIPROCIDAD_ES_EN_20260918/ARTICULO/ES/source/sections/electron.tex
\section{Estructura algebraica del electrón y escala de acción}
\label{x_reciprocidad:sec:electron}

La construcción electrónica recibe el estado aritmético con memoria de
trayectoria producido por APP--TRIT--TPK y sus coordenadas numéricas. En este punto del desarrollo,
\(\pi,\varphi,e\) y \(\alpha\) designan las coordenadas ya obtenidas de
la misma genealogía, con el orden constructivo expuesto anteriormente. No
son parámetros tomados de una tabla física. La sección electrónica es una
realización posterior de esa estructura discreta conjunta del continuo:
conserva la celda de origen, las dos hojas aritméticas, la orientación, el
acarreo y la memoria de retorno antes de producir un escalar.

Conviene separar tres objetos desde el comienzo. El primero es la sección
central y su historia de transporte; el segundo, la coordenada
adimensional que se obtiene de ella; el tercero, su realización en una
recta de energía o de masa. La igualdad entre dos coordenadas numéricas no
identifica estos tres niveles. En particular, el número de Euler \(e\),
la coordenada electrónica \(\varepsilon_e\) y una energía de reposo
\(E_e\) llevarán notaciones distintas.

El objetivo de esta sección es establecer la composición
\begin{equation}
 \varepsilon_e^{\beta}
 =\frac{\sqrt3}{4}
  \left[\exp\!\left(\frac{\varphi}{\pi^2}\right)-22\alpha^3\right]
  (1+15\Delta_4)
  R_{\mathrm{act}}^{\,80-54\alpha+6\Delta_4},
 \label{x_reciprocidad:eq:el-formula-completa}
\end{equation}
explicando por separado el origen del prefactor, los dos enteros de la
etapa basal, la normalización de \(\Delta_4\), la razón de acción y el
triple de retorno. La carta dimensional se abre solamente después de
esta composición. El resultado no consiste en hacer que un número
adimensional adquiera unidades por su mera semejanza con una magnitud
tabulada.

La composición reúne factores con funciones distintas. El cuadro
siguiente identifica su significado y remite a la construcción
correspondiente; las demostraciones se desarrollan en el orden de sus
dependencias.

\begin{table}[htbp]
\centering
\small
\setlength{\tabcolsep}{4pt}
\renewcommand{\arraystretch}{1.18}
\begin{tabularx}{\linewidth}{@{}p{0.20\linewidth}Xp{0.24\linewidth}@{}}
\toprule
Factor & Significado en la composición & Construcción y demostración\\
\midrule
\(\sqrt3/4\)
 & Norma del conmutador de los proyectores en el plano principal.
 & Ecuación~\eqref{x_reciprocidad:eq:el-prefactor}.\\
\(\exp(\varphi/\pi^2)\)
 & Carácter exponencial de la orientación electrónica del cociente regional.
 & Ecuación~\eqref{x_reciprocidad:eq:el-bisagra}; apartados~\ref{x_reciprocidad:sec:medida-retorno-hojas}
   y~\ref{x_reciprocidad:sec:operador-retorno-marcado}.\\
\(-22\alpha^3\)
 & Contribución del complemento regional en el retorno central.
 & Ecuaciones~\eqref{x_reciprocidad:eq:el-selector-enteros} y~\eqref{x_reciprocidad:eq:el-censo}.\\
\(1+15\Delta_4\)
 & Normalización torsional compuesta con la multiplicidad regional
   y trítica.
 & Ecuaciones~\eqref{x_reciprocidad:eq:el-selector-enteros}, \eqref{x_reciprocidad:eq:el-delta}
   y~\eqref{x_reciprocidad:eq:el-basal}.\\
\(R_{\mathrm{act}}^{\kappa_e}\)
 & Cociente de las secciones de acción, elevado al carácter determinado
   por los dos lectores del retorno electrónico.
 & Ecuaciones~\eqref{x_reciprocidad:eq:el-ratio}, \eqref{x_reciprocidad:eq:el-triple},
   \eqref{x_reciprocidad:eq:el-kappa} y~\eqref{x_reciprocidad:eq:el-beta}.\\
\bottomrule
\end{tabularx}
\caption{Factores de la composición electrónica y procedencia de sus reglas.}
\label{x_reciprocidad:tab:el-factores}
\end{table}


% BEGIN OWNER /Users/ruben/Documents/New project/output/EXTREMA_MEDIA_RAZON_RECIPROCIDAD_ES_EN_20260918/ARTICULO/ES/source/sections/electron_articulacion.tex
\paragraph{Identidad de la sección y composición de sus observables.}
\label{x_reciprocidad:el-ampl-articulacion}
La identidad estructural del electrón se introduce mediante la sección
central y sus operaciones de transporte. La inversión de la carta APP
selecciona la celda $(5,5)$ por su propiedad de punto fijo. Las dos hojas
aritméticas permiten estudiar las deformaciones que conservan la suma y
el residuo del producto. El resultado distingue las dos orientaciones
$(8,2)$ y $(2,8)$ y conserva la diferencia de una unidad de cociente.
Éste es el primer nivel de información: una celda, una involución, una
fibra conservada y un defecto entero mínimo. Ninguna evaluación de masa
es necesaria para enunciarlo o demostrarlo.

El paso a las proyecciones aritméticas introduce un segundo nivel. Las
fibras de suma y producto determinan esperanzas condicionales sobre el
mismo espacio finito. Su conmutador mide la dependencia respecto del
orden de esas dos proyecciones. La norma $\sqrt3/4$ resulta del espectro
de su emparejamiento; el bloque que alcanza esa norma conserva además
las relaciones cuadráticas demostradas a continuación. Por tanto, el
prefactor escalar y el álgebra local tienen una procedencia común, pero
contienen información diferente. La norma retiene la magnitud del
conmutador, mientras las matrices conservan su orientación y su ley de
composición.

La transferencia areal--volumétrica incorpora las coordenadas regionales
en un tercer nivel. El calendario de modos determina la matriz de
incidencia y su acción cuadrática. La marca de clausura produce la razón
$\pi/\varphi^2$; el intercambio de los dos argumentos proporciona la
orientación $\varphi/\pi^2$ utilizada en el carácter exponencial
electrónico. Las dos expresiones quedan relacionadas por una involución
explícita. Esa relación explica por qué una coordenada regional puede
aparecer en la ecuación electrónica con una potencia o en una posición
diferente de la que tenía en la primera lectura: la operación aplicada
es parte del dato matemático que se conserva.

El término cúbico de alfa pertenece al complemento regional de un retorno
marcado. La cuenta de sus posiciones produce $27-5=22$, mientras las
tres hojas de cada una de las cinco regiones producen $3\cdot5=15$.
Las demostraciones conservan el dominio de esos recuentos y la
orientación que determina sus signos. La fórmula basal compone,
por consiguiente, una norma de incompatibilidad, un carácter de
transferencia, un déficit de posiciones y una memoria torsional.
La escritura en una sola línea expresa la composición de esas
operaciones; su explicación exige reconstruirlas en ese orden.

La transición de la coordenada basal a la coordenada completa conserva
la misma sección central. El factor de retorno utiliza los registros
de la historia, y los dos lectores declarados recuperan el triple
$(80,54,6)$ sobre esa trayectoria. La multiplicación por
$R_{\mathrm{act}}^{80-54\alpha+6\Delta_4}$ modifica la evaluación,
pero mantiene identificados su origen y sus antecedentes. El valor
basal constituye una etapa intermedia de la composición; el valor
completo incluye también esa contribución del retorno. La comparación
física debe utilizar la etapa especificada por la fórmula que se evalúa.

La escala de acción interviene después con otra función. Una razón de
dos secciones de acción es adimensional y cancela su escala común.
La realización absoluta utiliza, en cambio, una base de acción, un
lector temporal y el mapa hacia energía o masa. Esta separación permite
conservar simultáneamente la construcción de la escala y la identidad
de la sección electrónica. El cuanto cronológico proporciona una
referencia de escala; la sección central determina el factor estructural
que se expresa respecto de ella. Las ecuaciones dimensionales de este
capítulo mantienen explícita esa composición.

La pertinencia del electrón en el argumento general procede de esta
convergencia de operaciones. La misma aritmética que produjo las
regiones determina ahora una celda central, sus deformaciones y sus
proyecciones. La constante de estructura fina, construida antes, actúa
en esa composición como coordenada de clausura. La información de
retorno determina la evaluación completa, y la realización dimensional
permite el contraste posterior. Explicar el electrón desde la estructura
consiste en conservar esta sucesión de objetos, mapas y valores en cada
paso de la demostración.

% END OWNER


\subsection{Centro, orientación y vacancia sobre una fibra conservada}

El soporte celular considerado es \(\{1,\ldots,9\}^2\), con ambas
evaluaciones APP, \(r+c\) y \(rc\). La inversión central es
\begin{equation}
 \rho(r,c)=(10-r,10-c).
 \label{x_reciprocidad:eq:el-inversion-central}
\end{equation}
El centro no se elige por una masa: se caracteriza por una propiedad del
soporte anterior a su lectura física.

\begin{proposition}[Unicidad del centro]
La inversión \(\rho\) tiene un único punto fijo, \((5,5)\).
\end{proposition}
\begin{proof}
Las ecuaciones \(r=10-r\) y \(c=10-c\) equivalen a
\(2r=2c=10\). Su única solución en la carta nonádica es \((5,5)\).
\end{proof}

La propiedad anterior identifica una celda, no una historia profunda
única. Una ruta añade la elección inicial de orientación y sus
prolongaciones compatibles. En la carta utilizada aquí la sección
orientada central comienza en \((5,5)\) con las dos orientaciones
positivas; se denotará por \(\Gamma_e\). El retorno de su proyección
celular no borra las orientaciones transportadas ni el registro de
acarreo de la sección con su memoria de trayectoria.

La vacancia se considera ahora, después del cierre de \(\alpha\), como
una operación local de esta sección. Su aritmética puede aislarse en un
lema finito; ese aislamiento no la convierte en un sustituto del
generador anterior de constantes.

\begin{definition}[Fibra central y defecto de cociente]
La fibra aditiva central es
\[
 \mathcal F_{10}
 =\{(5+t,5-t):t\in\mathbb Z,\ |t|\leq4\}.
\]
En la carta \(n=\rho_9(n)+9q_9(n)\), con
\(\rho_9(n)\in\{1,\ldots,9\}\), se llama defecto de cociente respecto
al centro a
\[
 v(t)=q_9(25)-q_9(25-t^2),
\]
cuando ambos productos poseen el mismo residuo nonádico.
\end{definition}

\begin{lemma}[Vacancia central mínima]
Los únicos puntos de \(\mathcal F_{10}\) que conservan el residuo
multiplicativo del centro son
\[
 (5,5),\qquad (8,2),\qquad (2,8).
\]
En los dos puntos no centrales el defecto de cociente es exactamente uno.
\end{lemma}
\begin{proof}
La suma es constantemente diez, mientras
\[
 (5+t)(5-t)=25-t^2.
\]
La conservación del residuo equivale a \(t^2\equiv0\pmod9\), es decir,
a \(3\mid t\). En el intervalo entero \([-4,4]\) las soluciones son
\(t=-3,0,3\). Para las dos no nulas,
\[
 25=7+9\cdot2,\qquad16=7+9\cdot1.
\]
Se conserva el residuo siete y se pierde una unidad de cociente. No hay
otro desplazamiento admisible no central en ese dominio, de modo que el
defecto positivo es mínimo. El intercambio \(t\mapsto-t\) permuta las
dos salidas y conserva el defecto.
\end{proof}

La vacancia no es aquí un cero numérico: es una diferencia entre dos
estados que tienen la misma lectura residual y distinto cociente. Si se
conservara únicamente el residuo, esa diferencia desaparecería. La
orientación distingue además los dos puntos de salida. Ni el defecto de
cociente ni esa orientación se identifican automáticamente con una carga
eléctrica o con una masa; su realización física requiere los lectores
correspondientes.


% BEGIN OWNER /Users/ruben/Documents/New project/output/EXTREMA_MEDIA_RAZON_RECIPROCIDAD_ES_EN_20260918/ARTICULO/ES/source/sections/centro_electronico_registros.tex
\subsubsection{Coordenadas centrales, orientación y registros regionales}
\label{x_reciprocidad:el-centro:registros}

La estructura central relaciona posición, evaluación y transporte. En la
celda \((5,5)\), las dos hojas de APP dan
\begin{equation}
 \mathcal A(5,5)=\bigl((1,1),(7,2)\bigr),\qquad
 5+5=1+9\cdot1,\qquad5\cdot5=7+9\cdot2.
 \label{x_reciprocidad:el-centro:app}
\end{equation}
Los dos unos de la hoja aditiva son, en su orden, residuo y cociente.
Su coincidencia numérica conserva dos funciones aritméticas distintas.
Sobre la fibra de suma \(10\), las dos vacancias orientadas verifican
\begin{equation}
 \mathcal A(8,2)=\mathcal A(2,8)=\bigl((1,1),(7,1)\bigr).
 \label{x_reciprocidad:el-centro:vacancias}
\end{equation}
Se conservan la hoja aditiva completa y el residuo multiplicativo;
la coordenada de cociente del producto disminuye exactamente en
\(1\). El transporte retiene, además, cuál de las dos posiciones
conjugadas se ha alcanzado.


% BEGIN OWNER /Users/ruben/Documents/New project/output/EXTREMA_MEDIA_RAZON_RECIPROCIDAD_ES_EN_20260918/ARTICULO/ES/source/figures/centro_y_vacancias.tex
\begin{figure}[H]
\centering
\begin{tikzpicture}[x=.67cm,y=.67cm,font=\footnotesize,>=Latex]
 \draw[step=1,black!12,line width=.3pt] (1,1) grid (9,9);
 \draw[->,black!60] (.7,1)--(9.5,1) node[right] {$i$};
 \draw[->,black!60] (1,.7)--(1,9.5) node[above] {$j$};
 \foreach \k in {1,...,9}{
  \node[below=2pt] at (\k,1) {$\k$};
  \node[left=2pt] at (1,\k) {$\k$};
 }
 \draw[black!65] (1,9)--(9,1);
 \draw[<->,line width=.9pt] (2,8)--(8,2);
 \foreach \x/\y in {2/8,5/5,8/2}{
  \filldraw[fill=white,line width=.8pt] (\x,\y) circle (3pt);
 }
 \fill (5,5) circle (1.5pt);
 \node[above right=3pt,fill=white,inner sep=1pt] at (2,8) {$(2,8)$};
 \node[above right=3pt,fill=white,inner sep=1pt] at (5,5) {$(5,5)$};
 \node[above right=3pt,fill=white,inner sep=1pt] at (8,2) {$(8,2)$};
 \node[anchor=west,align=left] at (10.5,7.6)
   {$i+j=10$\\$R^+=1,\quad q^+=1$};
 \node[anchor=west,align=left] at (10.5,5.4)
   {centro: $ij=25$\\$R^\times=7,\quad q^\times=2$};
 \node[anchor=west,align=left] at (10.5,3.1)
   {vacancias: $ij=16$\\$R^\times=7,\quad q^\times=1$};
\end{tikzpicture}
\caption{Centro electrónico y vacancias conjugadas en la carta APP.
La diagonal es la fibra aditiva \(\mathcal F_{10}\); el centro
es el punto fijo de \((i,j)\mapsto(10-i,10-j)\). Los dos
desplazamientos de módulo \(3\) conservan la suma y el residuo
producto y reducen en una unidad el cociente multiplicativo. Las
flechas representan las deformaciones sobre la fibra.}
\label{x_reciprocidad:el-centro:figura}
\end{figure}

% END OWNER


La simetría central caracteriza también su persistencia bajo cambios de
carta. Si una transformación admisible \(f\) del atlas entrelaza la
inversión, \(f\rho=\rho f\), entonces
\[
 \rho f(5,5)=f\rho(5,5)=f(5,5).
\]
La unicidad del punto fijo obliga a \(f(5,5)=(5,5)\). Ésta es la
propiedad interna que hace estable la localización central: toda
transformación de ese tipo conserva la celda antes de evaluar su
carácter electrónico (sección~\ref{x_reciprocidad:el-centro:registros}).

El bloque matricial central pertenece a la evaluación sobre las dos
posiciones adyacentes \(4,5\). Su descomposición, establecida en el
desarrollo local del núcleo, es
\[
 B_c=\begin{pmatrix}7&2\\2&7\end{pmatrix}=9P_++5P_-.
\]
Aquí \(9\) y \(5\) son autovalores de dos modos, mientras que
\((5,5)\) es una posición del atlas. Las concatenaciones de sus
entradas dan \(72+27=77+22=99\) y
\(77-22=55=54+1\). La primera igualdad conserva la suma total bajo
dos disposiciones; la segunda expresa el defecto entre las lecturas
diagonal y antidiagonal. Estas operaciones preceden a la composición
electrónica y conservan su procedencia matricial.

La orientación de partida es otro dato. En la representación del
retorno electrónico se fija \((h_0,v_0)=(1,1)\), es decir, avance
positivo en ambas coordenadas. La inversión simultánea cada \(27\)
pasos, leída en ventanas de \(9\), determina
\[
 \tau_e=(+1,+1,+1,-1,-1,-1,+1,+1,+1,-1,-1,-1).
\]
Los tramos \((1,1,1)\) pertenecen a esta sucesión de orientaciones:
cada tramo comprende tres ventanas del mismo sentido. La recurrencia
de la sección de lectores demuestra este registro en
\eqref{x_reciprocidad:eq:el-tau}, y sus dos evaluaciones recuperan el triple
\((80,54,6)\) que interviene en
\eqref{x_reciprocidad:eq:el-formula-completa}. La coordenada de posición, los
cocientes APP y la sucesión de signos quedan así conectados mediante
operaciones explícitas.

La firma \(555555\) pertenece a una lectura regional de otro
dominio. Para una emisión decimal \(U=(U_1,\ldots,U_6)\), escribimos
\(U_j=100d_{1j}+10d_{2j}+d_{3j}\), con tres cifras, incluidos los
ceros iniciales. El lector de firma multiplicativa es
\[
 \mathcal E_\times(U)
   =\bigl([d_{2j}-d_{1j}]_{10}\bigr)_{j=1}^{6}.
\]
En la región transversal
\(U_\pi^\perp=(501,614,498,169,272,272)\), las diferencias son
\((-5,-5,5,5,5,5)\). Por consiguiente,
\begin{equation}
 \mathcal E_\times(U_\pi^\perp)=(5,5,5,5,5,5),\qquad
 \Psi(U_\pi^\perp)=010211.
 \label{x_reciprocidad:el-centro:firma-transversal}
\end{equation}
La primera aplicación conserva una firma multiplicativa del catálogo;
la segunda conserva la palabra ternaria común a las cinco regiones.
La figura \ref{x_reciprocidad:pi-estrella:figura} muestra la posición transversal de
esta emisión y su multiplicidad \(432\).

\paragraph{Memoria necesaria para transportar la firma constante.}
La firma \(555555\) proporciona, además, un ejemplo verificable de
la información de trayectoria que exige TPK. Sea
\(\mathcal X_\times=D_9^2\times\{N,E,S,O\}\), con \(324\)
posiciones orientadas del canal producto. El avance \(P\) desplaza
una casilla en la dirección almacenada y la media vuelta \(H\)
invierte esa dirección. En esta representación, el lector
\(E_{\rm sig}\) suma por ventanas los exponentes discretos
\[
 \ell_9(1,2,4,8,7,5)=(0,1,2,3,4,5),\qquad
 \ell_9(3)=\ell_9(6)=\ell_9(9)=0.
\]
Durante cada una de las seis ventanas de nueve pasos, se evalúa
\(\ell_9(\rho_9(ij))\) antes de cada avance producto, seleccionado
por TRIT en los tiempos \(2,5,8\pmod9\). La suma de la ventana se
reduce módulo \(10\). Tras los pasos \(27\) y \(54\) se invierte
la dirección. Esta regla determina las seis componentes de
\(E_{\rm sig}\) desde la posición y orientación iniciales.

La enumeración de este dominio produce \(26\) firmas. La fibra de
\(555555\) contiene \(24\) posiciones orientadas; después de un
avance, sus imágenes se distribuyen en tres firmas, con ocho
representantes cada una:
\[
 555177,\qquad555717,\qquad555771.
\]
Tres testigos, en coordenadas de \(1\) a \(9\), permiten comprobar
la distinción directamente:
\[
\begin{array}{c|c|c}
x&E_{\rm sig}(x)&E_{\rm sig}(Px)\\\hline
(3,5,N)&555555&555177\\
(3,4,N)&555555&555717\\
(3,4,S)&555555&555771
\end{array}
\]
Por tanto, el valor común de la firma no determina su siguiente valor:
la clase de transporte dentro de esa fibra forma parte del dato
operativo. La partición estable mínima se calcula partiendo de la
equivalencia de firmas \(\sim_0\) y refinando mediante
\[
 x\sim_{n+1}y\quad\Longleftrightarrow\quad
 x\sim_n y,\quad Px\sim_n Py,\quad Hx\sim_n Hy.
\]
El procedimiento termina al ser finito el dominio. Toda congruencia
estable que refine \(\sim_0\) refina cada \(\sim_n\), por inducción;
su punto fijo es, por ello, la congruencia estable mínima requerida.
El censo da \(26\to28\to28\), con histograma
\(27\cdot8+108=324\). La única firma que se descompone es
\(555555\), en tres clases de ocho elementos. El certificado adjunto
reproduce la enumeración, las tres clases y la estabilidad por ambos
operadores (sección~\ref{x_reciprocidad:el-centro:registros}).

La multiplicidad \(24\) pertenece aquí a la fibra del canal producto;
la multiplicidad \(432\) pertenece a la emisión regional completa
\(U_\pi^\perp\). La primera prueba por qué la firma necesita memoria
de transporte; la segunda sitúa esa firma en el catálogo conjunto de
las dos hojas. Ambas lecturas participan de la misma arquitectura,
con dominios y operaciones especificados.

Esta distinción de dominios permite reunir las apariciones de la unidad
y del centro con sus respectivos significados:
\begin{center}
\small
\begin{tabularx}{\textwidth}{@{}lX@{}}
\toprule
Registro & Objeto y función \\
\midrule
\((5,5)\) & Celda fija de la inversión central del atlas APP.\\
\((R^+,q^+)=(1,1)\) & Residuo y cociente de la evaluación aditiva central.\\
\((h_0,v_0)=(1,1)\) & Orientación inicial de los dos desplazamientos.\\
\((+1,+1,+1)\) & Tres ventanas consecutivas de \(\tau_e\) con orientación positiva.\\
\(c=111111\) & Dato de la quinta frontera nonádica, sección \ref{x_reciprocidad:mono-recta:desarrollo}.\\
\(\mathcal E_\times=555555\) & Firma de la región transversal de \(\pi\), con seis posiciones.\\
\bottomrule
\end{tabularx}
\end{center}

La coordinación electrónica se obtiene componiendo estos niveles en su
orden de dependencia. APP determina el centro y las evaluaciones;
TRIT selecciona regímenes y orientaciones; TPK prolonga la trayectoria
y conserva sus registros. Las regiones previamente generadas aportan
los caracteres \(\pi,\varphi,e\); el registro dodecafásico aporta
\(\alpha\); los lectores electrónicos determinan los factores de
la ecuación \eqref{x_reciprocidad:eq:el-formula-completa}. Esta composición conserva
la identidad de la sección a lo largo de sus distintas
representaciones, hasta la evaluación dimensional y el contraste
metrológico posterior.

% END OWNER


\subsection{Norma del conmutador de las proyecciones aritméticas}

Se obtiene el prefactor electrónico sin utilizar \(\alpha\), la torsión
o una unidad dimensional. El dominio es la tercera realización de APP,
\(\mathcal A_3=\{1,\ldots,9\}^3\), provista de medida uniforme.
Definimos
\[
 \Sigma(x)=\rho_9(x_1+x_2+x_3),\qquad
 \Pi(x)=\rho_9(x_1x_2x_3).
\]
Las esperanzas condicionales \(P_\Sigma\) y \(P_\Pi\) son las
proyecciones ortogonales sobre las funciones constantes en cada fibra
aditiva y multiplicativa, respectivamente. Así, para una función \(f\),
\[
 (P_\Sigma f)(x)
 =\frac{1}{|\Sigma^{-1}(\Sigma(x))|}
   \sum_{y:\,\Sigma(y)=\Sigma(x)}f(y),
\]
y análogamente para \(P_\Pi\). Las dos proyecciones actúan sobre el
mismo espacio; no representan dos soportes aritméticos independientes.

\begin{lemma}[Dos proyecciones y ángulos principales]
Si \(s\in[0,1]\) es un valor singular del emparejamiento entre bases
ortonormales de los rangos de dos proyecciones \(P,Q\), su bloque no
trivial contribuye a \(\|[P,Q]\|\) con
\(s\sqrt{1-s^2}\).
\end{lemma}
\begin{proof}
En el plano principal correspondiente se pueden escribir
\[
 P=\begin{pmatrix}1&0\\0&0\end{pmatrix},\qquad
 Q=\begin{pmatrix}
 s^2&s\sqrt{1-s^2}\\
 s\sqrt{1-s^2}&1-s^2
 \end{pmatrix}.
\]
Su conmutador es
\[
 [P,Q]=s\sqrt{1-s^2}
 \begin{pmatrix}0&1\\-1&0\end{pmatrix}.
\]
La matriz final es ortogonal, por lo que su norma es uno. Los bloques
comunes u ortogonales tienen conmutador nulo. La descomposición ortogonal
en planos principales da la afirmación y reduce la norma total al máximo
de estas contribuciones.
\end{proof}

\begin{proposition}[Norma exacta de incompatibilidad]
Sobre \(\ell^2(\mathcal A_3)\),
\begin{equation}
 \|[P_\Sigma,P_\Pi]\|=\frac{\sqrt3}{4}.
 \label{x_reciprocidad:eq:el-prefactor}
\end{equation}
\end{proposition}
\begin{proof}
Sea \(N_{ab}=|\Sigma^{-1}(a)\cap\Pi^{-1}(b)|\). La enumeración entera
de los tres símbolos proporciona
\[
 N=9\begin{pmatrix}
 0&1&1&0&1&1&0&1&4\\
 1&0&1&1&0&1&1&0&4\\
 1&0&2&0&1&2&0&0&3\\
 0&1&1&0&1&1&0&1&4\\
 1&0&1&1&0&1&1&0&4\\
 0&0&2&1&0&2&0&1&3\\
 0&1&1&0&1&1&0&1&4\\
 1&0&1&1&0&1&1&0&4\\
 0&1&2&0&0&2&1&0&3
 \end{pmatrix}.
\]
Cada fila suma \(81\); las sumas de columna son
\[
 (m_1,\ldots,m_9)=(36,36,108,36,36,108,36,36,297).
\]
Las indicatrices normalizadas de las fibras tienen matriz de
emparejamiento \(C_{ab}=N_{ab}/\sqrt{81m_b}\). Por consiguiente,
\[
 M=CC^{\mathsf T},\qquad
 M_{ac}=\sum_{b=1}^9\frac{N_{ab}N_{cb}}{81m_b}
\]
es una matriz racional cuyo espectro consta de los cuadrados de los
valores singulares requeridos.

La multiplicación exacta muestra que \((1,\ldots,1)\),
\((1,-1,0)^3\) y \((1,1,-2)^3\) son autovectores de autovalores
\(1\), \(1/4\) y \(7/132\), respectivamente. El subespacio de
vectores soportados en \(\{3,6,9\}\) cuya suma es cero tiene dimensión
dos y autovalor \(1/18\). Los vectores de suma cero soportados en
\(\{1,4,7\}\) o en \(\{2,5,8\}\) forman un núcleo de dimensión
cuatro. Estos subespacios son independientes y sus dimensiones suman
nueve. Queda, pues, probado el espectro completo
\[
 \operatorname{spec}(M)=
 \left\{1,\frac14,\frac7{132},\frac1{18},\frac1{18},0,0,0,0\right\}.
\]
El máximo de \(\lambda(1-\lambda)\) en este conjunto es \(3/16\),
alcanzado en \(\lambda=1/4\). El lema anterior da
\(\sqrt{3/16}=\sqrt3/4\).
\end{proof}

El modo \((1,-1,0)^3\) es precisamente la distribución de valores del
selector trítico de fase en la carta nonádica. La norma escalar retiene
una medida de incompatibilidad, pero no conserva por sí sola la
orientación de ese modo ni reconstruye los dos proyectores. Esa memoria
permanece en el estado de ruta que acompaña la lectura escalar.

\begin{proposition}[Álgebra del bloque electrónico local]
En el plano principal que alcanza \eqref{x_reciprocidad:eq:el-prefactor}, sean
\[
 P=\begin{pmatrix}1&0\\0&0\end{pmatrix},\qquad
 Q=\begin{pmatrix}1/4&\sqrt3/4\\\sqrt3/4&3/4\end{pmatrix},\qquad
 S=2P-I,\qquad J=\frac4{\sqrt3}[P,Q].
\]
Entonces \(S^2=I\), \(J^2=-I\), \(SJ=-JS\), y el álgebra real
generada por \(S,J\) es \(M_2(\mathbb R)\), realización de
\(\mathrm{Cl}_{1,1}\). Además,
\[
 n_\pm=\frac{S\pm J}{2},\qquad
 n_\pm^2=0,\qquad n_+n_-+n_-n_+=I.
\]
\end{proposition}
\begin{proof}
Se tiene \(S=\operatorname{diag}(1,-1)\) y
\(J=\bigl(\begin{smallmatrix}0&1\\-1&0\end{smallmatrix}\bigr)\).
La multiplicación verifica las tres relaciones. Las matrices
\(I,S,J,SJ\) son linealmente independientes y forman una base de
\(M_2(\mathbb R)\), lo que prueba la afirmación sobre el álgebra.
Finalmente, \((S\pm J)^2=S^2+J^2\pm(SJ+JS)=0\), y la suma de los
productos cruzados es \((S^2-J^2)/2=I\).
\end{proof}

En este bloque conviven una dirección elíptica, una hiperbólica y dos
direcciones nulas parabólicas. El resultado describe una representación
matricial local de la aritmética de las dos hojas. No identifica toda APP
con \(M_2(\mathbb R)\), ni con un cuadrado latino cuántico; tales
afirmaciones exigirían otros dominios y otras relaciones. Tampoco la
exponencial pertenece exclusivamente al régimen parabólico: puede
aplicarse a cualquiera de los generadores de este bloque.


% BEGIN OWNER /Users/ruben/Documents/New project/output/EXTREMA_MEDIA_RAZON_RECIPROCIDAD_ES_EN_20260918/ARTICULO/ES/source/sections/helicidad_electron.tex
% Adición focal propuesta; no sustituye el álgebra electrónica existente.
% Inserción: después de la proposición «Álgebra del bloque electrónico local»
% y su párrafo explicativo, antes de «Transferencia entre hojas...».
% Procedencia y alcance: LOCALIZADORES_Y_ALCANCE.md, en esta carpeta.
\subsection{Representación espinorial del bloque electrónico y helicidad}
\label{x_reciprocidad:sec:el-representacion-espinorial}

La estructura electrónica dispone de una representación que conserva más
información que su evaluación escalar. El plano principal de los dos
proyectores de APP tridimensional, seleccionado por el modo trítico
\((1,-1,0)^3\), ya ha proporcionado los operadores \(S\) y \(J\).
El primero es una involución simétrica; el segundo es el conmutador
normalizado y conserva la orientación de las hojas. Su composición permite
construir explícitamente una representación de rotaciones y separar sus
dos componentes de helicidad. Esta construcción actúa sobre el bloque
electrónico; el registro TPK de la ruta y su memoria permanecen como datos
de procedencia del bloque.

Sea \(E_e\) el plano real principal, dotado del producto escalar heredado
de las esperanzas condicionales, y sea
\(\mathscr S_e=E_e\otimes_{\mathbb R}\mathbb C\) su complejificación.
Los operadores ya obtenidos satisfacen
\[
 S^*=S,\qquad J^*=-J,\qquad
 S^2=I,\qquad J^2=-I,\qquad SJ=-JS.
\]
La complejificación extiende el cuerpo de representación. Conserva los
operadores reales y permite distinguir la unidad escalar \(i\) del
endomorfismo \(J\), aunque ambos tengan cuadrado \(-1\) en sus
respectivos dominios.

\begin{proposition}[Terna hermítica inducida por las dos hojas]
\label{x_reciprocidad:prop:el-terna-hermitica}
Los operadores
\begin{equation}
 \sigma_1=SJ,\qquad \sigma_2=-iJ,\qquad \sigma_3=S
 \label{x_reciprocidad:eq:el-terna-hermitica}
\end{equation}
son hermíticos y de traza nula. Satisfacen
\begin{equation}
 \sigma_a\sigma_b
 =\delta_{ab}I+i\sum_{c=1}^3\epsilon_{abc}\sigma_c,
 \qquad
 \frac12\operatorname{tr}(\sigma_a\sigma_b)=\delta_{ab},
 \label{x_reciprocidad:eq:el-algebra-terna}
\end{equation}
donde \(\epsilon_{123}=1\). Por tanto constituyen una base ortonormal
del espacio real \(V_e\) de endomorfismos hermíticos de traza nula de
\(\mathscr S_e\), para el producto
\(\langle X,Y\rangle=\frac12\operatorname{tr}(XY)\).
\end{proposition}
\begin{proof}
Las relaciones anteriores dan \((SJ)^*=(-J)S=SJ\) y
\((-iJ)^*=-iJ\). Los cuadrados de las tres matrices son \(I\).
Además,
\[
 (SJ)(-iJ)=iS,\qquad (-iJ)S=iSJ,\qquad S(SJ)=J=i(-iJ).
\]
Al invertir el orden cambia el signo. Esto prueba la primera identidad
de \eqref{x_reciprocidad:eq:el-algebra-terna}. En la base ortonormal del plano principal,
las matrices son
\[
 \sigma_1=\begin{pmatrix}0&1\\1&0\end{pmatrix},\quad
 \sigma_2=\begin{pmatrix}0&-i\\i&0\end{pmatrix},\quad
 \sigma_3=\begin{pmatrix}1&0\\0&-1\end{pmatrix}.
\]
Su traza es nula, y tomar trazas en la identidad de producto prueba la
ortonormalidad. Todo endomorfismo hermítico de traza nula tiene la forma
\(\bigl(\begin{smallmatrix}z&x-iy\\x+iy&-z\end{smallmatrix}\bigr)\),
con \(x,y,z\) reales, por lo que pertenece a su envolvente real.
\end{proof}

Esta terna recibe, en la representación matricial usual, el nombre de
matrices de Pauli. Su empleo aquí es posterior a la construcción de
\(P,Q,S,J\): se obtiene por los productos
\eqref{x_reciprocidad:eq:el-terna-hermitica}. El espacio de direcciones de esta
representación es la esfera unidad de \(V_e\), cuya forma positiva
acaba de construirse mediante la traza.

\begin{proposition}[Generadores de rotación y descomposición direccional]
\label{x_reciprocidad:prop:el-helicidad}
Defínanse \(L_a=\sigma_a/2\). Entonces
\begin{equation}
 [L_a,L_b]=i\sum_c\epsilon_{abc}L_c,
 \qquad \sum_{a=1}^3L_a^2=\frac34 I.
 \label{x_reciprocidad:eq:el-generadores-spin}
\end{equation}
Para \(n=(n_1,n_2,n_3)\) con \(\sum_an_a^2=1\), sean
\[
 H_e(n)=\sum_an_a\sigma_a,\qquad
 h_e(n)=\frac12H_e(n),\qquad
 \Pi_\pm(n)=\frac12\bigl(I\pm H_e(n)\bigr).
\]
La fórmula lineal \(H_e(v)=\sum_av_a\sigma_a\) se utiliza también
para \(v\) de norma arbitraria.
Se tiene
\begin{equation}
 \mathscr S_e=\operatorname{im}\Pi_+(n)
             \mathbin{\oplus}\operatorname{im}\Pi_-(n),
 \qquad
 h_e(n)\Pi_\pm(n)=\pm\frac12\Pi_\pm(n).
 \label{x_reciprocidad:eq:el-descomposicion-helicidad}
\end{equation}
Los dos sumandos son líneas complejas ortogonales. El transporte de
rotación alrededor de \(n\) es
\begin{equation}
 U_n(\theta)=\exp\!\left(-\frac{i\theta}{2}H_e(n)\right)
 =\cos\frac\theta2\,I-i\sin\frac\theta2\,H_e(n),
 \qquad
 U_n(2\pi)=-I,\quad U_n(4\pi)=I.
 \label{x_reciprocidad:eq:el-retorno-spin}
\end{equation}
\end{proposition}
\begin{proof}
La identidad de producto de la terna da los conmutadores de
\eqref{x_reciprocidad:eq:el-generadores-spin}; sus cuadrados suman \(3I/4\).
Los términos antisimétricos se cancelan al multiplicar
\(H_e(n)\) por sí mismo, de modo que \(H_e(n)^2=I\).
Se deduce
\[
 \Pi_\pm(n)^2=\Pi_\pm(n),\qquad
 \Pi_+(n)\Pi_-(n)=0,\qquad \Pi_+(n)+\Pi_-(n)=I.
\]
Son hermíticos y tienen traza \(1\); por ello su rango es \(1\).
Multiplicar por \(H_e(n)/2\) prueba
\eqref{x_reciprocidad:eq:el-descomposicion-helicidad}. La serie de la exponencial se
separa en potencias pares e impares gracias a \(H_e(n)^2=I\), lo que
produce \eqref{x_reciprocidad:eq:el-retorno-spin}. También prueba directamente la
unitariedad y la ley \(U_n(\theta+\psi)=U_n(\theta)U_n(\psi)\).

La normalización angular se verifica sobre \(V_e\). Para todo
\(v\in\mathbb R^3\), la misma identidad de producto da
\[
 U_n(\theta)H_e(v)U_n(\theta)^*
 =H_e\!\left(v\cos\theta+(n\mathbin{\times}v)\sin\theta
       +n(n\mathbin{\cdot}v)(1-\cos\theta)\right).
\]
La acción inducida es una rotación de ángulo \(\theta\) y eje \(n\).
El factor \(1/2\) en el generador espinorial es, por tanto, el que
corresponde a esa parametrización angular. Una rotación completa restaura
la dirección en \(V_e\) y cambia el signo del vector de
\(\mathscr S_e\); dos rotaciones completas restauran ambos.
\end{proof}

La representación es irreducible: un subespacio complejo invariante por
las tres matrices sería invariante por toda \(M_2(\mathbb C)\), el
álgebra que generan. Sus únicos subespacios invariantes son \(0\) y
\(\mathscr S_e\). Las relaciones
\eqref{x_reciprocidad:eq:el-generadores-spin} identifican así la representación de
espín \(1/2\). Si \(n\) es la dirección de propagación en una
realización cinemática, \(h_e(n)\) es su helicidad adimensional, con
valores \(\pm1/2\). La dirección cinemática se declara al efectuar esa
realización; los dos proyectores están definidos para cada dirección
antes de seleccionar un estado de propagación particular.

\paragraph{Conservación de la estructura electrónica y tipos de inversión.}
La inversión de dirección tiene una ley exacta:
\[
 H_e(-n)=-H_e(n),\qquad \Pi_\pm(-n)=\Pi_\mp(n).
\]
Esta operación intercambia las etiquetas de helicidad relativas a una
dirección. La conjugación de la carga actúa, en cambio, sobre la firma
orientada de la ruta material; la inversión celular actúa sobre las
coordenadas del atlas. Sus dominios y acciones permanecen distintos.
La unicidad de la celda \((5,5)\) es compatible con las dos líneas de
\eqref{x_reciprocidad:eq:el-descomposicion-helicidad}: una celda del atlas no es un
vector de la representación espinorial que porta.

La evaluación electrónica \(\varepsilon_e^\beta\) conserva su
fórmula y sus lectores. Sobre \(\mathscr S_e\), su realización escalar
es \(\varepsilon_e^\beta I\), que conmuta con los generadores y con
\(\Pi_\pm(n)\). La descomposición de helicidad añade la información
rotacional sin modificar la evaluación ni convertir el espín en un
coeficiente multiplicativo de la masa.

La quiralidad corresponde a la graduación de la representación relativista
que se utilice: una involución \(\Gamma_{\rm ch}\) determina sus
proyectores \((I\mp\Gamma_{\rm ch})/2\). La helicidad de
\eqref{x_reciprocidad:eq:el-descomposicion-helicidad} depende además de \(n\).
Por su parte, la holonomía de una banda real de Möbius describe el cambio
de signo de una línea al recorrer un lazo de su base. La ley de esa línea,
el retorno espinorial \eqref{x_reciprocidad:eq:el-retorno-spin} y la holonomía nonádica
de fase con avance de memoria son leyes de transporte sobre objetos
especificados. Su comparación conserva esos objetos: la forma helicoidal
de una trayectoria no sustituye la representación que acaba de
construirse, ni fija por sí misma la quiralidad o la carga.

% END OWNER


\subsection{Transferencia entre hojas y coeficientes de la etapa basal}

El calendario trítico primitivo visita los modos
\((\Sigma,\Pi,\Pi)\). Al contar los pares consecutivos, incluido el
cierre del bloque, se obtienen una transición \(\Sigma\to\Pi\), una
\(\Pi\to\Pi\) y una \(\Pi\to\Sigma\). Su incidencia primitiva es
\[
 F_{\mathrm{av}}=\begin{pmatrix}0&1\\1&1\end{pmatrix}.
\]
Esta realización matricial se obtiene después del calendario. No se
emplea para reemplazar el lector coinductivo que ya generó \(\varphi\).
Su polinomio característico es \(X^2-X-1\), y el reconocimiento de las
dos raíces como \(\varphi\) y \(-\varphi^{-1}\) describe la acción de
esa salida en la transferencia modal. La acción cuadrática tiene autovalores
\(\varphi^2,-1,\varphi^{-2}\); el último es el modo estable positivo.
La marca regional \(\pi\), aplicada una sola vez, da
\(\pi/\varphi^2\).

El espacio de trayectorias, su medida y el operador de clausura están
construidos en los apartados~\ref{x_reciprocidad:sec:medida-retorno-hojas}
y~\ref{x_reciprocidad:sec:operador-retorno-marcado}. La fórmula
\eqref{x_reciprocidad:eq:retorno-marcado-finito} conserva su procedencia como cociente
de retornos; la aplicación electrónica compone después esta lectura
con la involución siguiente. La firma regional \(555555\) acompaña
a la región de clausura, cuyo transporte conserva el estado con memoria;
el centro \((5,5)\) fija la celda de origen de la sección electrónica
orientada. Ambas informaciones convergen en la composición y mantienen
sus dominios respectivos.

\begin{lemma}[Intercambio de las dos orientaciones]
Sean \(\mathcal R(x,y)=x/y^2\) y
\(\mathscr S_{\pi\varphi}(x,y)=(y,x)\), para \(x,y>0\).
Entonces
\begin{equation}
 \begin{aligned}
 \mathscr S_{\pi\varphi}^{\,2}&=I,\qquad
 \mathcal R(\pi,\varphi)=\frac{\pi}{\varphi^2},\\
 (\mathcal R\circ\mathscr S_{\pi\varphi})(\pi,\varphi)
 &=\frac{\varphi}{\pi^2}
 =\frac1{\varphi^3(\pi/\varphi^2)^2}.
 \end{aligned}
 \label{x_reciprocidad:eq:el-bisagra}
\end{equation}
\end{lemma}
\begin{proof}
El intercambio aplicado dos veces es la identidad. Además,
\(\varphi^3(\pi/\varphi^2)^2=\pi^2/\varphi\); tomar su inversa prueba
la última igualdad.
\end{proof}

El lector electrónico utiliza la orientación \(\varphi/\pi^2\),
seguida del carácter exponencial. La identidad anterior explica el
intercambio, pero no sustituye la regla que elige ese carácter ni produce
por sí sola los enteros de la corrección. Estos poseen otra procedencia.

El retorno central compatible con la fase tiene longitud mínima
veintisiete. Con el registro de fase y hoja, su espacio finito es
\[
 \mathcal K_e=\mathbb Z/9\mathbb Z\times\mathbb F_3.
\]
Las cinco regiones de \(\pi\) son identidades genealógicas distintas,
no cinco copias intercambiables de una expansión visible. En el calibre
regional fijado por la construcción, su inyección en el retorno tiene
los lugares
\begin{equation}
 (4,0),\quad(5,0),\quad(6,0),\quad(6,1),\quad(6,2).
 \label{x_reciprocidad:eq:el-cinco-lugares}
\end{equation}
La orientación transversal, el retorno mutado y las tres regiones
positivas permanecen separados por ese registro. Los cambios admisibles
de calibre transportan las etiquetas conjuntamente; no vuelven a
seleccionar los lugares a partir de un valor electrónico.

\begin{proposition}[Déficit y memoria retenida]
Fijados el retorno central y la inyección regional anterior, el déficit
orientado de frontera y la memoria torsional son
\begin{equation}
 n_e=-|\mathcal K_e\setminus\iota(\mathscr R_\pi)|=-22,
 \qquad
 m_e^{\mathrm{tor}}=|\mathscr R_\pi\times\mathbb F_3|=15.
 \label{x_reciprocidad:eq:el-selector-enteros}
\end{equation}
\end{proposition}
\begin{proof}
Los cinco lugares de \eqref{x_reciprocidad:eq:el-cinco-lugares} son distintos. Por
tanto, \(|\mathcal K_e|=9\cdot3=27\) y el complemento de la imagen
tiene cardinal \(27-5=22\). La convención de orientación registra como
negativo el déficit retirado de la frontera. La memoria retenida posee
tres hojas por región, de modo que su cardinal positivo es
\(5\cdot3=15\).
\end{proof}

La misma pareja posee el control entero de la carta APP de la familia:
\begin{equation}
 396=18+72+108+198,
 \qquad \frac{396}{18}=22,
 \qquad \frac{270}{18}=15=3\cdot5=\binom62.
 \label{x_reciprocidad:eq:el-censo}
\end{equation}
El peso \(18\) pertenece al núcleo central \(2\times2\), y \(270\)
al canal global declarado. Estas igualdades controlan la normalización
del recuento; el argumento de retorno fija además los signos. El
cociente de los anillos alternos,
\(396/(18+108)=22/7\), es racional y no se identifica con \(\pi\).
Tampoco \(\binom62\) significa \(6/2\): cuenta los pares no ordenados
de un conjunto de seis elementos.

La evaluación del déficit en la tercera realización simétrica emplea
\(\alpha^3\). Esta potencia es el producto de tres copias del escalar
de cierre ya obtenido; no afirma una regla universal que haga
corresponder cualquier dimensión con la misma potencia de \(\alpha\).
El dominio y el lector fijados son esenciales para la composición.

\subsection{La torsión global y la coordenada basal}

\begin{definition}[Normalización torsional electrónica]
Se distinguen el defecto reducido \(d_4\) y el lector completo
\(\Delta_4\):
\begin{equation}
 d_4=\frac{\pi-e\log\pi}{270},\qquad
 \Delta_4=d_4\left(1+\frac{\pi}{729}\right).
 \label{x_reciprocidad:eq:el-delta}
\end{equation}
El denominador \(270\) corresponde al canal del censo, mientras el
factor \(1+\pi/729\) retiene el acoplamiento con la celda cúbica de
\(729\) lugares. La lectura electrónica utiliza \(\Delta_4\) tanto en
su etapa basal como en el exponente de retorno.
\end{definition}

Esta definición actúa sobre \(\pi\) y \(e\) previamente generados. El
logaritmo es una operación posterior en su realización real, no una
entrada retrospectiva de valores convencionales. La torsión es global:
no se determina independientemente para cada sección mediante una masa
observada.

\begin{lemma}[Las dos normalizaciones no son intercambiables]
Se cumple
\begin{equation}
 \Delta_4-d_4=\frac{\pi}{729}d_4.
 \label{x_reciprocidad:eq:el-delta-diferencia}
\end{equation}
En particular, \(d_4>0\) y \(\Delta_4>d_4\) para las coordenadas
consideradas.
\end{lemma}
\begin{proof}
La identidad es distributiva. Para la positividad, la función
\(f(x)=x-e\log x\), sobre \(x>0\), tiene derivada \(1-e/x\) y
segunda derivada \(e/x^2>0\). Su mínimo único está en \(x=e\), donde
vale cero. Como \(\pi\ne e\), se tiene \(f(\pi)>0\). Ambos factores
de \eqref{x_reciprocidad:eq:el-delta} son positivos y \(1+\pi/729>1\).
\end{proof}

\begin{remark}[Conservación de la normalización]
Suprimir el segundo factor de \eqref{x_reciprocidad:eq:el-delta} cambia el lector; no
es una simplificación algebraica. Si se introduce \(d_4\) como
coordenada auxiliar, el mismo objeto debe escribirse
\(\Delta_4=(1+\pi/729)d_4\) en todas sus apariciones. Mantener los
enteros \(15\) y \(6\) y sustituir sólo \(\Delta_4\) por \(d_4\)
produce otra fórmula.
\end{remark}

\begin{definition}[Coordenada electrónica basal]
La composición de incompatibilidad, transferencia modal, déficit y memoria retenida
es
\begin{equation}
 \varepsilon_e^{(0)}
 =\frac{\sqrt3}{4}
  \left[\exp\!\left(\frac{\varphi}{\pi^2}\right)-22\alpha^3\right]
  (1+15\Delta_4).
 \label{x_reciprocidad:eq:el-basal}
\end{equation}
\end{definition}

La ecuación reúne tres operaciones de naturaleza distinta. La norma
\(\sqrt3/4\) cuantifica la incompatibilidad de las proyecciones de
las hojas; su origen es un ángulo principal de las fibras de APP.
La exponencial actúa sobre la orientación \(\varphi/\pi^2\) del
lector de transferencia. La resta cúbica registra las \(22\) posiciones
complementarias del retorno marcado, y el factor torsional conserva las
\(15\) componentes correspondientes a tres hojas por región. La
composición conserva estos dominios aunque su evaluación final sea escalar.

La presencia de \(\pi\), \(\varphi\) y \(e\) tiene así significados
operativos. El número de Euler interviene por la exponencial y por
la normalización logarítmica de \(\Delta_4\); \(\varphi\) y \(\pi\)
intervienen en la transferencia orientada. La constante \(\alpha\)
actúa después de su propia construcción mediante el término cúbico
de clausura. El retorno cronológico aún debe operar sobre esta coordenada
basal. La separación de etapas permite explicar qué cambia al modificar
una orientación, una normalización o un registro, conservando la estructura
central que identifica la sección electrónica.

Los paréntesis de esta ecuación son parte de su contenido. El término
\(-22\alpha^3\) se resta \emph{fuera} de la exponencial; la memoria
\(1+15\Delta_4\) multiplica todo el corchete. Introducir el déficit en
el exponente o aplicarle una torsión distinta alteraría la composición.

\begin{proposition}[Positividad basal]
Para \(0<\alpha<10^{-2}\), \(\varphi>0\) y la normalización
\eqref{x_reciprocidad:eq:el-delta}, se tiene \(\varepsilon_e^{(0)}>0\).
\end{proposition}
\begin{proof}
La exponencial es mayor que uno, y
\(22\alpha^3<22\cdot10^{-6}<1\). Por ello el corchete es positivo.
También \(1+15\Delta_4>1\) y \(\sqrt3/4>0\).
\end{proof}

La evaluación posterior de \eqref{x_reciprocidad:eq:el-basal} comienza por
\[
 \varepsilon_e^{(0)}=0.5109989224880660725\ldots.
\]
Esta cifra es todavía una coordenada sin dimensión y no constituye la
última etapa: falta aplicar la memoria de retorno de la propia ruta.

\subsection{Escala determinantal de acción y valoración decimal}

La escala absoluta y el corrector relativo se originan en objetos
distintos. El primero utiliza el bloque local del registro dodecafásico
de paso tres,
\begin{equation}
 H_4=\begin{pmatrix}
 1&1&1&1\\1&1&-1&-1\\1&-1&1&-1\\1&-1&-1&1
 \end{pmatrix},\qquad
 G_H=\mathbb Z^4/H_4\mathbb Z^4.
 \label{x_reciprocidad:eq:el-h4}
\end{equation}
La operación utiliza un cociente local. Los tres bloques del registro
completo no se sustituyen por una tensorización que multiplicara de
nuevo el número de hojas de este lector.

\begin{lemma}[Número de hojas del cociente local]
La forma normal de Smith de \(H_4\) es
\(\operatorname{diag}(1,2,2,4)\). En consecuencia,
\[
 G_H\simeq\mathbb Z/2\mathbb Z\oplus\mathbb Z/2\mathbb Z
 \oplus\mathbb Z/4\mathbb Z,\qquad |G_H|=16.
\]
\end{lemma}
\begin{proof}
El máximo común divisor de las entradas es uno. Los menores de orden dos
son pares y hay uno de valor \(-2\), de modo que su máximo común divisor
es dos. Como \(H_4H_4^{\mathsf T}=4I\) y \(\det H_4=-16\), los
cofactores son \(\pm4\); el tercer divisor determinantal es cuatro.
El cuarto es \(16\). Los cocientes sucesivos de los divisores
\(1,2,4,16\) son \(1,2,2,4\). Su producto da el orden del cociente.
\end{proof}

\begin{definition}[Lector determinantal de acción]
La norma multiplicativa de la sección escalar \(\alpha\), constante
sobre las hojas de \(G_H\), y una aplicación de la autoescala
\(\varphi\) producen
\begin{equation}
 \Sigma_H=\varphi\,\mathrm N_{G_H}(\alpha),\qquad
 \mathrm N_{G_H}(\alpha)=\prod_{g\in G_H}\alpha=\alpha^{16}.
 \label{x_reciprocidad:eq:el-lector-determinantal}
\end{equation}
La carta decimal asociada determina
\(\operatorname{ord}_{10}(\Sigma_H)=\lfloor\log_{10}\Sigma_H\rfloor\).
\end{definition}

La potencia dieciséis procede así del número de hojas del cociente, una
vez fijado este lector local. El factor \(\varphi\) actúa una vez sobre
la base de la recta de acción. El cálculo de Smith, por sí solo, no
seleccionaría cualquier funcional posible sobre el cociente: la norma
multiplicativa y la autoescala forman parte de la operación declarada.

\begin{theorem}[Década de la sección de acción]
El lector \eqref{x_reciprocidad:eq:el-lector-determinantal}, aplicado a las coordenadas
HMT ya generadas, satisface
\begin{equation}
 10^{-34}<\Sigma_H<10^{-33},\qquad
 \operatorname{ord}_{10}(\Sigma_H)=-34.
 \label{x_reciprocidad:eq:el-decada}
\end{equation}
\end{theorem}
\begin{proof}
Bastan los intervalos racionales, mucho más gruesos que los cilindros
disponibles,
\[
 \frac{1618}{1000}<\varphi<\frac{1619}{1000},\qquad
 \frac{7297}{10^6}<\alpha<\frac{7298}{10^6}.
\]
La función \((x,a)\mapsto xa^{16}\) es creciente en ambos argumentos
positivos. Las dos desigualdades enteras
\[
 1618\cdot7297^{16}>10^{65},\qquad
 1619\cdot7298^{16}<10^{66}
\]
implican
\[
 10^{-34}
 <\frac{1618}{1000}\left(\frac{7297}{10^6}\right)^{16}
 <\Sigma_H
 <\frac{1619}{1000}\left(\frac{7298}{10^6}\right)^{16}
 <10^{-33}.
\]
La definición del orden decimal concluye la prueba. No interviene una
unidad SI ni una masa de referencia en estas desigualdades.
\end{proof}

La sección normalizada comienza por
\(\Sigma_H=1.0462438381\ldots\times10^{-34}\). Su mantisa no se
identifica con la coordenada de acción anterior o posterior al retorno.
El lector determinantal fija el orden decimal; las dos coordenadas de
retorno se obtienen mediante los lectores siguientes.


\subsection{Secciones de acción y cociente adimensional de retorno}

Sea \(\mathcal L_S\) una recta orientada de acción y \(\mathcal U_S\)
su base positiva. La composición electrónica utiliza dos secciones de
esta recta y su cociente adimensional. Su construcción requiere la
década determinantal ya obtenida, el carácter de acción de quinto orden
y la contribución regional del retorno.

\subsubsection{Carácter de acción y procedencia de sus coeficientes}
\label{x_reciprocidad:sec:revision-planck}

La constante de Planck reducida, \(\hbar\), es el objeto físico que
motiva la realización de la sección de acción angular. Su posición en
la dependencia electrónica debe especificarse en dos niveles. La norma
determinantal \(\Sigma_H\) determina la década de la escala. El carácter
\(H_5\) y el retorno regional determinan después las coordenadas de las
secciones de acción. Esta separación permite exponer el resultado
necesario para el electrón sin anticipar los desarrollos completos de
otros sectores de la mecánica dimensional.

La procedencia del carácter de quinto orden puede explicitarse mediante
invariantes ya fijados en el soporte. En la carta central se considera
\[
 B_c=\begin{pmatrix}7&2\\2&7\end{pmatrix},
 \qquad \lambda_+=9,\quad\lambda_-=5.
\]
El calendario tiene longitud \(12\), la órbita del paso \(3\) tiene
longitud \(4\), y el cociente censal pertinente es
\(104976/1944=54\). Los coeficientes del carácter satisfacen
\begin{equation}
 \frac9{16}=\frac{\lambda_+}{4^2},\quad
 \frac59=\frac{\lambda_-}{\lambda_+},\quad
 \frac7{48}=\frac{\operatorname{tr}(B_c)/2}{4\cdot12},\quad
 \frac1{54}=\frac{1944}{104976}.
 \label{x_reciprocidad:eq:revision-h5-coeficientes}
\end{equation}
La asignación de estos invariantes a los grados \(2,3,4,5\), con la
alternancia de signos de \eqref{x_reciprocidad:eq:el-h5}, pertenece a la definición del
carácter analítico utilizado. Explicitarla es esencial: el conjunto de
cardinales, considerado sin el carácter y sin el orden de sus grados,
admitiría otras expresiones. El contenido constructivo reside en la
composición del soporte, la orientación y la regla analítica declarada.

Sobre esta base, el carácter de quinto orden y el término regional de
retorno tienen las expresiones
\begin{align}
 H_5={}&\frac12\sqrt{\frac{1000\alpha}{\varphi}}-\alpha
   +\frac9{16}\alpha^2-\frac59\alpha^3
   +\frac7{48}\alpha^4-\frac1{54}\alpha^5,
   \label{x_reciprocidad:eq:el-h5}\\
 D_A={}&\exp\!\left(-\frac{100\pi\alpha}{9}\right),
   \label{x_reciprocidad:eq:el-da}\\
 \mathcal C_\pi(\alpha)={}&169\alpha^6+
       \frac{D_A\alpha^7}{1-D_A\alpha},
   \label{x_reciprocidad:eq:el-cpi}\\
 \eta_{\mathrm{ret}}={}&H_5-\frac{90}{\pi}\mathcal C_\pi(\alpha).
   \label{x_reciprocidad:eq:el-eta}
\end{align}
Estas son lecturas posteriores del cierre dodecafásico y de su
incidencia regional. El entero \(169\) se conserva como el selector
regional ya construido, y los coeficientes de \(H_5\) como los de su
carácter local; no se reemplazan por una interpolación de la constante de
Planck. Las fórmulas explicitan exactamente qué lecturas se necesitan
para la etapa electrónica. No es necesario introducir aquí otros
funcionales de la teoría de masas.

Para \(0<\alpha<1\), se tiene \(0<D_A<1\), de modo que
\(D_A\alpha<1\). El término racional de \eqref{x_reciprocidad:eq:el-cpi} conserva la
suma completa de la cola geométrica
\begin{equation}
 \frac{D_A\alpha^7}{1-D_A\alpha}
 =\sum_{n=0}^{\infty}D_A^{n+1}\alpha^{n+7}.
 \label{x_reciprocidad:eq:el-cola-regional}
\end{equation}
No es una truncación que pueda modificarse silenciosamente al cambiar la
precisión. Tras el término de índice \(N\), la cola restante es
\(D_A^{N+2}\alpha^{N+8}/(1-D_A\alpha)\), lo que permite controlar la
evaluación sin acudir a una magnitud objetivo.

\begin{definition}[Secciones de acción y lector de retorno]
Las secciones anterior y posterior al retorno son
\begin{equation}
 \hbar_{\mathrm{pre}}=H_5\,10^{-34}\mathcal U_S,
 \qquad
 \hbar_{\mathrm{ret}}=\eta_{\mathrm{ret}}\,10^{-34}\mathcal U_S.
 \label{x_reciprocidad:eq:el-secciones-accion}
\end{equation}
Siempre que \(\eta_{\mathrm{ret}}>0\), se definen
\begin{equation}
 \delta_{\mathrm{ret}}=H_5-\eta_{\mathrm{ret}},\qquad
 R_{\mathrm{act}}=\frac{\hbar_{\mathrm{pre}}}{\hbar_{\mathrm{ret}}}
 =\frac{H_5}{\eta_{\mathrm{ret}}}.
 \label{x_reciprocidad:eq:el-ratio}
\end{equation}
\end{definition}

\begin{proposition}[Positividad e invariancia de la razón]
En los intervalos usados en la prueba de \eqref{x_reciprocidad:eq:el-decada},
\(0<\eta_{\mathrm{ret}}<H_5\), y por tanto
\(R_{\mathrm{act}}>1\). La razón no cambia al sustituir la base
\(\mathcal U_S\) por cualquier otra base positiva común.
\end{proposition}
\begin{proof}
Los términos de \(\mathcal C_\pi\) son positivos. Para una cota
elemental basta usar \(\alpha<0.008\), \(\varphi<1.619\),
\(\alpha>0.007\), \(D_A<1\) y \(\pi>3\). La raíz en
\eqref{x_reciprocidad:eq:el-h5} es mayor que uno, mientras la suma de los términos
negativos es menor que \(0.008001\); en particular, \(H_5>0.99\).
Asimismo,
\[
 0<\frac{90}{\pi}\mathcal C_\pi
 <30\left(169(0.008)^6+
           \frac{(0.008)^7}{1-0.008}\right)<10^{-6}.
\]
Así \(\eta_{\mathrm{ret}}>0.989999>0\) y
\(\eta_{\mathrm{ret}}<H_5\). Finalmente, el cambio de base divide
ambas coordenadas de acción por el mismo factor positivo, que cancela en
su cociente.
\end{proof}

La evaluación posterior proporciona
\begin{align*}
 H_5&=1.0545718176461544696\ldots,\\
 \eta_{\mathrm{ret}}&=1.0545718169150390611\ldots,\\
 R_{\mathrm{act}}&=1.0000000006932817630\ldots.
\end{align*}
La diferencia y el cociente son dos lecturas de las mismas secciones,
una aditiva y otra multiplicativa. No representan dos constantes de
Planck independientes. El paso de acción angular a acción por vuelta
completa se expresa, en cada sección, mediante
\(h=2\pi\hbar\). La posterior carta
\(\mathcal U_S\mapsto\mathrm{J\,s}\) asigna una unidad, pero no
retroactúa sobre el orden decimal ya demostrado.

El retorno regional puede escribirse también en una carta angular.
Esta reparametrización no es una dependencia adicional indispensable:
\eqref{x_reciprocidad:eq:el-h5}--\eqref{x_reciprocidad:eq:el-eta} suministran directamente
\(H_5\), \(\eta_{\mathrm{ret}}\) y \(R_{\mathrm{act}}\). Si se
introducen ángulo y contraángulo, ambos deben transportarse en la misma
carta; no se mezclan grados y radianes dentro de una diferencia.


% BEGIN OWNER /Users/ruben/Documents/New project/output/EXTREMA_MEDIA_RAZON_RECIPROCIDAD_ES_EN_20260918/ARTICULO/ES/source/sections/revision_planck.tex
\subsubsection{Continuación regional y ecuación de renovación}

El selector \(169\) determina el impulso de grado seis. La prolongación
posterior se organiza por la transferencia \(D_A\) y una normalización
unitaria de la línea determinantal. Estas reglas admiten una formulación
en series formales antes de sustituir \(z=\alpha\), con \(D_A\) ya
construido. Escribamos
\[
 \mathscr C(z)=169z^6+R(z),\qquad R(z)\in z^7\mathbb R[[z]].
\]
La concatenación de una prolongación aumenta el grado en una unidad y
multiplica su peso por \(D_A\). La primera prolongación estricta tiene
peso \(D_A\). Por tanto, su ecuación de renovación es
\begin{equation}
 R(z)=D_Az^7+D_AzR(z).
 \label{x_reciprocidad:eq:revision-renovacion}
\end{equation}

\begin{proposition}[Unicidad formal de la continuación normalizada]
La ecuación \eqref{x_reciprocidad:eq:revision-renovacion} determina una única serie y
\[
 \mathscr C(z)=169z^6+\frac{D_Az^7}{1-D_Az}
             =169z^6+\sum_{n=7}^{\infty}D_A^{n-6}z^n.
\]
La realización es analítica en \(|D_Az|<1\).
\end{proposition}
\begin{proof}
El elemento \(1-D_Az\) tiene término independiente uno y es invertible
en el anillo de series formales. Resolver la ecuación produce la serie
exhibida. Equivalentemente, el coeficiente de grado siete es \(D_A\)
y cada coeficiente posterior es \(D_A\) veces el anterior. La serie
geométrica converge absolutamente en el disco indicado.
\end{proof}

La normalización de la primera prolongación es un dato operativo
indispensable. Si se conservaran únicamente el impulso de grado seis y
la razón de concatenación, las series
\[
 169z^6+\lambda\frac{D_Az^7}{1-D_Az}
\]
seguirían compartiendo ambos datos para cualquier \(\lambda\).
La regla unitaria fija \(\lambda=1\) antes de efectuar la evaluación.
Así queda localizada la contribución de cada regla a la unicidad del
lector, y la expresión racional conserva la continuación completa.

\subsubsection{Constante de Planck reducida y retorno de la sección de acción}

En la carta dimensional de \eqref{x_reciprocidad:eq:el-secciones-accion}, la sección
anterior al retorno es la construcción del modelo destinada al contraste
con la constante de Planck reducida:
\[
 \hbar_{\mathrm{pre}}=H_5\,10^{-34}\mathcal U_S.
\]
La sección posterior conserva, además, la compensación regional
\(90\mathscr C(\alpha)/\pi\). Ambas pertenecen a la misma recta de
acción y su cociente \(R_{\mathrm{act}}\) mide el efecto multiplicativo
del retorno. El paso a \(h\) mediante \(2\pi\) utiliza la relación
entre la parametrización angular y una vuelta completa.

Esta construcción interviene efectivamente en
\(R_{\mathrm{act}}^{80-54\alpha+6\Delta_4}\). El electrón conserva
su coordenada central y su trayectoria; la escala cronológica no
reemplaza esa estructura por un cuanto indiferenciado. Tampoco se
multiplica de nuevo por \(10^{-34}\) una masa ya expresada en MeV:
la década pertenece a la sección de acción, cancela en su cociente
adimensional y reaparece únicamente donde actúa la realización
dimensional correspondiente.

El contraste numérico de \(H_5\), de la sección retornada y de
\(h/(2\pi)\) se presenta en la sección~\ref{x_reciprocidad:sec:revision-planck-contraste}.
Así se distingue la denominación del objeto físico, la ecuación que el
modelo le asigna y la precisión con la que dicha ecuación lo reproduce.

% END OWNER


\subsection{Dos lectores del retorno electrónico}

La memoria electrónica se representa en una palabra trítica de longitud
\(108\) y en una palabra firmada de doce fases. Para que el cálculo del
exponente sea reproducible, explicitamos la proyección cronológica
utilizada. Sean inicialmente
\((r_0,c_0)=(5,5)\) y \((h_0,v_0)=(1,1)\). En cada paso el selector
repite \((+1,-1,0)\). Con orientación fija dentro de cada bloque de
veintisiete pasos, la regla visible es
\[
 \begin{array}{c|c|c}
 \text{modo}&\text{entero leído}&\text{transporte celular}\\\hline
 +1&r+c&c\mapsto1+(c-1+h)\bmod9\\
 -1&rc&r\mapsto1+(r-1+v)\bmod9\\
 0&9&\text{sin desplazamiento}
 \end{array}
\]
El símbolo emitido es \(\rho_9(\text{entero leído})\bmod3\).
Después de los pasos \(27,54,81\) se invierten simultáneamente
\(h\) y \(v\). Esta regla calcula la proyección observable; la
actualización del estado conserva adicionalmente residuo, cociente,
acarreo, hoja, frontera y memoria. No se identifica el registro de
\(108\) símbolos con el estado completo.

En cada bloque completo de tres pasos se ha avanzado una vez en cada
coordenada; hacen falta nueve de estos bloques para volver al centro en
la misma fase. La longitud mínima de este retorno registrado es, por
tanto, \(27\). Si se olvida la fase, la celda ya se alcanza tras el
paso \(26\), antes del umbral neutro que completa el bloque; ése no es
el retorno tipado que cuenta \(\mathcal K_e\). La lectura explícita
de los dos primeros bloques de retorno da
\begin{equation}
 a=100000220,\qquad b=120200000,\qquad
 \gamma_e=(a^3b^3)^2.
 \label{x_reciprocidad:eq:el-palabra}
\end{equation}
Aquí las potencias denotan concatenación de palabras. En particular,
\(\gamma_{e,t+54}=\gamma_{e,t}\). En el registro de nueve pasos por
fase, las orientaciones centrales dan
\begin{equation}
 \tau_e=(1,1,1,-1,-1,-1,1,1,1,-1,-1,-1).
 \label{x_reciprocidad:eq:el-tau}
\end{equation}
La periodicidad de estas representaciones no equivale al reinicio de la
historia que las produce.

\begin{definition}[Lector de discrepancias cíclicas]
Para una palabra \(\gamma\in\mathbb F_3^{108}\), se define
\[
 D_k(\gamma)=\#\{t\in\mathbb Z/108\mathbb Z:
                   \gamma_t\ne\gamma_{t+k}\}.
\]
En el dominio de palabras de período divisor de \(54\), el lector
entero es
\begin{equation}
 K_D(\gamma)=\left(D_{27}+D_{36},D_1,\frac{D_4-D_3}{2}\right).
 \label{x_reciprocidad:eq:el-kd}
\end{equation}
\end{definition}

\begin{lemma}[Integridad y evaluación central]
Los \(D_k\) del dominio anterior son pares. Para
\eqref{x_reciprocidad:eq:el-palabra},
\[
 \begin{array}{c|rrrrr}
 k&1&3&4&27&36\\\hline
 D_k&54&56&68&48&32
 \end{array}
\]
y \(K_D(\gamma_e)=(80,54,6)\).
\end{lemma}
\begin{proof}
La aplicación \(t\mapsto t+54\) empareja sin puntos fijos los índices
de discrepancia y conserva su condición; cada cardinal es par. Para la
palabra explícita \(a^3b^3\) de longitud \(54\), el recuento de
discrepancias con cada desplazamiento es, en el mismo orden,
\(27,28,34,24,16\). Duplicar la palabra duplica esos cinco cardinales.
Finalmente,
\[
 (D_{27}+D_{36},D_1,(D_4-D_3)/2)
 =(48+32,54,(68-56)/2)=(80,54,6).
\]
\end{proof}

Los desplazamientos \(27\) y \(36\) son las elevaciones
cronológicas de los canales de \(90\) y \(120\) grados en una carta
donde nueve pasos corresponden a treinta grados. La lectura angular
explica esa representación, pero los enteros se computan sobre la
palabra, sin introducir una masa ni un valor real buscado.

\begin{definition}[Lector de ventanas dodecafásicas]
Para \(\tau\in\{-1,0,1\}^{12}\), con índices cíclicos, sean
\begin{align*}
 C_k(\tau)&=\sum_{r=0}^{11}\tau_r\tau_{r+k},\\
 A_\ell(\tau)&=\sum_{r=0}^{11}
           \left(\sum_{j=0}^{\ell-1}\tau_{r+j}\right)^2,\\
 P_6(\tau)&=\sum_{r=0}^{5}\tau_r\tau_{r+6}.
\end{align*}
Se define
\begin{equation}
 \Omega(\tau)=4A_3(\tau)-3A_4(\tau),\qquad
 K_\Omega(\tau)=(\Omega(\tau),54,P_6(\tau)).
 \label{x_reciprocidad:eq:el-komega}
\end{equation}
El \(54\) de este lector registra la semivuelta cronológica declarada.
No se confunde con el salto local de valor cuatro de APP.
\end{definition}

\begin{lemma}[Contraste primitivo y forma de correlación]
Entre los contrastes lineales enteros \(uA_3+vA_4\) que se anulan
siempre que \(A_3=3a\) y \(A_4=4a\), la pareja primitiva, con
orientación \(90-120\), es \((u,v)=(4,-3)\). Además,
\begin{equation}
 \Omega=-2C_1-4C_2-6C_3.
 \label{x_reciprocidad:eq:el-omega-corr}
\end{equation}
\end{lemma}
\begin{proof}
La anulación para todo \(a\) exige \(3u+4v=0\). Por coprimalidad,
las soluciones son \((u,v)=n(4,-3)\), y la orientación fija el signo
del generador primitivo. Al desarrollar los cuadrados,
\[
 A_\ell=\ell C_0+2\sum_{k=1}^{\ell-1}(\ell-k)C_k.
\]
Por tanto,
\(A_3=3C_0+4C_1+2C_2\) y
\(A_4=4C_0+6C_1+4C_2+2C_3\). Su combinación \(4A_3-3A_4\)
cancela \(C_0\) y produce \eqref{x_reciprocidad:eq:el-omega-corr}.
\end{proof}

La unicidad demostrada pertenece a esa clase de contrastes lineales
enteros; no excluye por sí sola todo funcional no lineal de las ventanas.
La especificación de la clase funcional evita convertir una identidad
de dos canales en una afirmación de selección universal.

\begin{theorem}[Triple de retorno de la sección electrónica]
Los dos lectores aplicados a los registros de \(\Gamma_e\)
coinciden y dan
\begin{equation}
 K_D(\gamma_e)=K_\Omega(\tau_e)=(80,54,6).
 \label{x_reciprocidad:eq:el-triple}
\end{equation}
Su evaluación en las coordenadas ya generadas es
\begin{equation}
 \kappa_e=80-54\alpha+6\Delta_4.
 \label{x_reciprocidad:eq:el-kappa}
\end{equation}
\end{theorem}
\begin{proof}
El primer lector se calculó arriba. Para \(\tau_e\), el producto
antipodal es siempre uno, y así \(P_6=6\). El recuento cíclico da
\(C_1=4\), \(C_2=-4\), \(C_3=-12\); luego
\[
 \Omega=-2\cdot4-4(-4)-6(-12)=80.
\]
Equivalentemente, \(A_3=44\), \(A_4=32\) y
\(4\cdot44-3\cdot32=80\). Por tanto ambos triples son
\((80,54,6)\). Evaluar el funcional
\((A,B,C)\mapsto A-\alpha B+\Delta_4C\) prueba la última fórmula.
\end{proof}

Esta coincidencia corresponde a la sección central. No identifica los
dos lectores en todos los dominios de rutas: uno compara símbolos
tríticos y el otro compara ventanas de un registro firmado. Tampoco el
\(80\) se obtiene del período de una palabra asignada al número de
Euler. Su origen en esta composición es el triple de la ruta
electrónica.

\subsection{Composición electrónica completa y normalización}

\begin{definition}[Retorno multiplicativo]
La acción de la memoria de retorno sobre la coordenada basal es
\begin{equation}
 \varepsilon_e^\beta
 =\varepsilon_e^{(0)}\exp\!\left(\kappa_e\log R_{\mathrm{act}}\right).
 \label{x_reciprocidad:eq:el-beta}
\end{equation}
\end{definition}

\begin{theorem}[Factorización electrónica positiva]
Con los lectores, la orientación y la normalización declarados,
\eqref{x_reciprocidad:eq:el-beta} coincide con \eqref{x_reciprocidad:eq:el-formula-completa}. Todos
sus factores son adimensionales y la coordenada resultante es positiva.
\end{theorem}
\begin{proof}
La norma de incompatibilidad es \eqref{x_reciprocidad:eq:el-prefactor}; la transferencia modal,
el déficit y la memoria retenida componen \eqref{x_reciprocidad:eq:el-basal}. El
triple \eqref{x_reciprocidad:eq:el-triple} produce \eqref{x_reciprocidad:eq:el-kappa}, y la razón de
las dos secciones de acción es \eqref{x_reciprocidad:eq:el-ratio}. Sustituyendo esas
tres identidades en \eqref{x_reciprocidad:eq:el-beta} se obtiene exactamente
\eqref{x_reciprocidad:eq:el-formula-completa}. La etapa basal es positiva y la
exponencial de un número real también. El factor común de dimensión de
acción ya canceló al formar \(R_{\mathrm{act}}\).
\end{proof}

La evaluación de esta composición comienza por
\begin{equation}
 \varepsilon_e^\beta
 =0.5109989506900008086082571648\ldots.
 \label{x_reciprocidad:eq:el-beta-decimal}
\end{equation}
Las cifras son una evaluación posterior de la fórmula, no datos usados
para seleccionar \(22\), \(15\) o \((80,54,6)\). La representación
numérica del basal mediante un prefijo finito introduce únicamente su
error de truncación; no modifica la fórmula exacta que lo define.

La incidencia de una normalización torsional distinta puede localizarse
sin compararla con una medida. Escribamos
\(B=\frac{\sqrt3}{4}[\exp(\varphi/\pi^2)-22\alpha^3]\),
\(R=R_{\mathrm{act}}\) y
\[
 \mathcal E(d)=B(1+15d)R^{80-54\alpha+6d}.
\]
Para los dos defectos de \eqref{x_reciprocidad:eq:el-delta},
\begin{equation}
 \frac{\mathcal E(\Delta_4)}{\mathcal E(d_4)}
 =\frac{1+15\Delta_4}{1+15d_4}
   R^{6(\Delta_4-d_4)}>1.
 \label{x_reciprocidad:eq:el-normalizacion-beta}
\end{equation}
La primera desigualdad torsional y \(R>1\) prueban que no son la misma
coordenada. Si sólo se cambia el defecto del exponente, manteniendo el
basal completo, la razón de las dos salidas es
\(R^{6(\Delta_4-d_4)}\), también distinta de uno. Una diferencia muy
pequeña sigue siendo una diferencia de fórmula; no puede atribuirse a
una simple conversión de unidades.

\subsection{Realización dimensional: energía, masa y escala cronológica}

La coordenada electrónica actúa finalmente sobre una recta dimensional.
Sea \(\mathcal L_E\) una recta de energía con base positiva
\(\mathcal U_E\). La realización es
\begin{equation}
 E_e^{(0)}=\varepsilon_e^{(0)}\mathcal U_E,\qquad
 E_e^\beta=\varepsilon_e^\beta\mathcal U_E,
 \qquad m_e^\beta=\frac{E_e^\beta}{c_{\mathrm{int}}^2}.
 \label{x_reciprocidad:eq:el-dimensional}
\end{equation}
La última expresión pertenece a la recta de masa cuando
\(c_{\mathrm{int}}\) lleva la dimensión de velocidad en la realización
declarada. En la carta comparativa \(\mathcal U_E=1\,\mathrm{MeV}\),
\eqref{x_reciprocidad:eq:el-beta-decimal} expresa una energía en MeV; la masa
correspondiente se expresa en \(\mathrm{MeV}/c^2\). Es frecuente
abreviar ambas magnitudes tomando \(c=1\), pero esa abreviatura no
debe ocultar el mapa dimensional.

La escala de acción también admite una realización acción--reloj. Si
\(t_0\) es una sección temporal positiva, el reloj de \(108\) pasos
define
\begin{equation}
 \omega_0=\frac{2\pi}{108t_0},\qquad
 E_{\mathrm{clk}}=\hbar_{\mathrm{ret}}\omega_0,\qquad
 m_{\mathrm{clk}}=E_{\mathrm{clk}}c_{\mathrm{int}}^{-2}.
 \label{x_reciprocidad:eq:el-cuanto-reloj}
\end{equation}
La acción por frecuencia tiene dimensión de energía. Este cuanto
cronológico es una sección dimensional común; no es, por definición,
la sección electrónica. La palabra de ruta y sus factores no se
obtienen de la uniformidad del reloj.

\begin{proposition}[Cancelación de la escala común]
El factor \(10^{-34}\mathcal U_S\) de
\eqref{x_reciprocidad:eq:el-secciones-accion} no forma parte de \(\kappa_e\). En
una misma recta de energía,
\begin{equation}
 \frac{E_e^\beta}{E_e^{(0)}}=R_{\mathrm{act}}^{\kappa_e}.
 \label{x_reciprocidad:eq:el-razon-energia}
\end{equation}
El resultado es independiente de la elección de la base dimensional
común.
\end{proposition}
\begin{proof}
Al formar \(\hbar_{\mathrm{pre}}/\hbar_{\mathrm{ret}}\), el factor
\(10^{-34}\mathcal U_S\) cancela exactamente. Al dividir las dos
expresiones de energía de \eqref{x_reciprocidad:eq:el-dimensional}, cancela además
\(\mathcal U_E\), y \eqref{x_reciprocidad:eq:el-beta} da la razón indicada.
\end{proof}

Por tanto, una coordenada ya expresada en MeV no se multiplica otra vez
por \(10^{-34}\). Si se elige \(E_{\mathrm{clk}}\) como base de
\(\mathcal L_E\), deben transportarse las coordenadas por el factor
de cambio de base; conservar sin cambio la cifra de la carta de un MeV
sería identificar dos bases sin demostrarlo. Éste es el punto preciso
en que la escala de acción entra en la realización absoluta y deja de
intervenir en el corrector relativo.

Existe asimismo una realización lineal que hace visible la separación
entre reloj y genealogía. En
\(\ell^2(\mathbb Z/108\mathbb Z;\mathbb R)\), el desplazamiento
cíclico \(S\) tiene espacio fijo
\[
 \ker(S-I)=\mathbb R n_0,\qquad
 n_0=\frac1{\sqrt{108}}\sum_{t=0}^{107}\delta_t.
\]
Si \(g_{\Gamma_e}\) es el símbolo de la ruta en la linealización libre
de las genealogías y \(\mathcal U_M\) una base de la recta de masa,
queda definido el mapa de rango uno
\begin{equation}
 \mathfrak C_e:
 \mathbb R n_0\otimes\mathbb R g_{\Gamma_e}\longrightarrow\mathcal L_M,
 \qquad
 \mathfrak C_e(n_0\otimes g_{\Gamma_e})
 =\varepsilon_e^\beta\mathcal U_M.
 \label{x_reciprocidad:eq:el-rango-uno}
\end{equation}
El mapa reúne la factorización anterior una vez declaradas sus bases.
El vector uniforme \(n_0\) no selecciona \(\varepsilon_e^\beta\):
esa coordenada procede de APP, del retorno orientado y de sus lectores.
Así se distingue una realización algebraica del modo nulo de una
identificación del electrón con el cuanto del reloj.

La conclusión interna es la factorización positiva de una sección
central, con sus operaciones, memoria y normalizaciones explícitas. La
asignación a un electrón físico añade la representación de carga, el
espín, el acoplamiento y la carta de unidades correspondiente; no se
deduce exclusivamente de una coincidencia decimal. El contraste
metrológico se presenta después de la prolongación excepcional y no altera
las entradas de ninguna de las construcciones aquí realizadas.


% BEGIN OWNER /Users/ruben/Documents/New project/output/EXTREMA_MEDIA_RAZON_RECIPROCIDAD_ES_EN_20260918/ARTICULO/ES/source/sections/revision_electron_argumento.tex
\subsubsection{Dependencia estructural de la coordenada electrónica}

La ecuación electrónica expresa una composición jerárquica. Su primer
factor, \(\sqrt3/4\), pertenece a la estructura matricial central y fija
una normalización geométrica. El exponente \(\varphi/\pi^2\) utiliza las
coordenadas regionales ya generadas en la orientación indicada del lector
areal--volumétrico. El término \(22\alpha^3\) incorpora el acoplamiento
dodecafásico en la corrección cúbica. El factor \(1+15\Delta_4\) conserva
la normalización torsional global. Finalmente, el cociente de las secciones
de acción actúa con el exponente determinado por los lectores de retorno.
Cada operación tiene así un antecedente distinto dentro de la misma
construcción.

La articulación resulta más precisa al considerar sus dependencias.
La norma central puede calcularse antes de evaluar las coordenadas
regionales. La exponencial y la corrección cúbica necesitan esas
coordenadas y el cierre de \(\alpha\). La memoria multiplicativa necesita
además las secciones de acción y la trayectoria de \(108\) posiciones.
Por ello, conocer la expresión basal no equivale a haber evaluado la
coordenada completa. La sucesión de operaciones no agrega una masa
externa al punto \((5,5)\); desarrolla los lectores atribuidos a su
sección persistente.

La distinción entre escala y estructura proporciona también una
interpretación de la comparación física. La masa es la realización
dimensional de la coordenada completa, mientras el centro, la involución
y la vacancia describen las relaciones que la producen. Un cambio común
de unidad modifica las coordenadas dimensionales, pero deja inalteradas
las razones adimensionales y las identidades de transporte demostradas.
Éste es el sentido en que el electrón conserva su carácter estructural
al expresarse respecto de una escala cronológica.

% END OWNER



% BEGIN OWNER /Users/ruben/Documents/New project/output/EXTREMA_MEDIA_RAZON_RECIPROCIDAD_ES_EN_20260918/ARTICULO/ES/source/sections/realizacion_electromagnetica.tex
\newpage
\subsection{Transporte lorentziano y representación electromagnética}

La sección electrónica proporciona una estructura y una coordenada
antes de su aplicación a una interacción. La realización lorentziana
de la incidencia cúbica determina, con la convención temporal declarada,
el espacio de formas en el que puede expresarse esa interacción.
Para formularla se fija una realización local orientada \((\mathcal M,g)\)
de dimensión cuatro, con fibras tangentes identificadas con el modelo
\(Q_\square\). Se especifican además una conexión \(\nabla\) con
\(\nabla g=0\), una uno-forma potencial \(A\in\Omega^1(\mathcal M)\),
su curvatura \(F=dA\) y una representación de carga de la sección.
Estos datos son los argumentos de la representación física que sigue.

La dualidad de Hodge satisface \(\star^2=-I\) sobre las dos-formas.
La separación eléctrica y magnética corresponde a una elección temporal
en la descomposición \(1+3\); \(F\) y \(\star F\) conservan su
relación dentro del mismo espacio de dos-formas. La neutralidad de una
representación de carga significa \(q=0\), mientras la sección puede
conservar orientación, fase y transporte interno. La ausencia de
acoplamiento eléctrico en esa representación tiene un contenido distinto
de la ausencia de información de fase.

En esta realización dimensional, sean \(m_e>0\), \(q\) la carga,
\(\gamma(s)\) una trayectoria parametrizada y \(u=\dot\gamma\)
su cuadrivelocidad. La fuerza de Lorentz se expresa por
\begin{equation}
 m_e\nabla_u u^\mu=qF^\mu{}_{\nu}u^\nu.
 \label{x_reciprocidad:eq:lorentz-electron}
\end{equation}
La antisimetría de \(F\) implica
\(F_{\mu\nu}u^\mu u^\nu=0\). Por consiguiente,
\[
 \frac{d}{ds}g(u,u)=2g(\nabla_u u,u)=0.
\]
Se conserva la normalización de la trayectoria siempre que \(\nabla\)
sea métrica y se cumpla \eqref{x_reciprocidad:eq:lorentz-electron}. La masa utilizada
es la realización dimensional de la sección electrónica precedente;
el parámetro de acoplamiento se identifica en la convención electromagnética
adoptada. La deducción concreta del potencial y de una configuración
de campo constituye un problema adicional con condiciones de frontera
propias. El resultado aquí expuesto determina el dominio, los datos
de la representación y la conservación que su ecuación implica.

La inversión de rutas del TPK y la orientación de una línea electrónica
ofrecen entonces una comparación precisa: cada una tiene su involución,
su transporte y su representación. El antecedente de Wheeler y Feynman
sitúa históricamente la pregunta; la equivalencia de dos realizaciones
se decide mediante sus aplicaciones y condiciones de compatibilidad.
Esta disposición permite mantener unido el origen aritmético del electrón
con la formulación de su interacción, haciendo explícita la función
de cada etapa.

% END OWNER

